\documentclass[11pt]{article}
\usepackage{cornell-notes}
\usepackage{mathtools}
\usepackage{needspace}
\usepackage{ragged2e}
\tcbuselibrary{most}
\geometry{letterpaper,margin=0.72in,headheight=15pt}
\setCornellCollection{Cloud Architecture Cornell Notes}
\setCornellUnitType{Chapter}
\setCornellUnitNumber{4}
\setCornellUnitTitle{Fundamental Concepts and Models}

\definecolor{Navy}{HTML}{17324D}
\definecolor{Teal}{HTML}{147D78}
\definecolor{Sky}{HTML}{E8F2F7}
\definecolor{CueGray}{HTML}{EEF2F4}
\definecolor{Warm}{HTML}{FFF3D6}
\definecolor{Risk}{HTML}{953B35}
\definecolor{Green}{HTML}{276749}
\definecolor{RuleGray}{HTML}{AAB8C2}
\definecolor{Text}{HTML}{25313A}

\hypersetup{colorlinks=true,linkcolor=Navy,urlcolor=Teal,citecolor=Navy,pdfborder={0 0 0}}
\color{Text}
\setlength{\parindent}{0pt}
\setlength{\parskip}{4pt}
\setlength{\emergencystretch}{3em}
\setlength{\LTpre}{3pt}
\setlength{\LTpost}{5pt}
\renewcommand{\arraystretch}{1.16}
\setlist{nosep,leftmargin=1.45em}
\hfuzz=0.2pt
\hbadness=2500

\titleformat{\section}{\Large\bfseries\color{Navy}\RaggedRight}{}{0pt}{}[\color{Teal}\titlerule]
\titleformat{\subsection}{\large\bfseries\color{Teal}\RaggedRight}{}{0pt}{}
\titlespacing*{\section}{0pt}{1.35ex plus .4ex}{.7ex}
\titlespacing*{\subsection}{0pt}{1.1ex plus .3ex}{.45ex}

\pagestyle{fancy}
\fancyhf{}
\fancyhead[L]{\small\color{Navy}\textbf{Cloud Architecture Cornell Notes}}
\fancyhead[R]{\small\color{Teal}\ChapterMark}
\fancyfoot[L]{\footnotesize\color{Navy}Concepts $\bullet$ Mechanisms $\bullet$ Evidence $\bullet$ Decisions}
\fancyfoot[R]{\footnotesize\color{Navy}\thepage\ of \pageref*{LastPage}}
\renewcommand{\headrulewidth}{0.5pt}
\renewcommand{\footrulewidth}{0.35pt}

\newtcolorbox{orientationbox}{enhanced,breakable,colback=Sky,colframe=Teal,boxrule=.75pt,arc=2mm,title={Chapter Orientation},fonttitle=\bfseries,left=2mm,right=2mm,top=1.5mm,bottom=1.5mm}
\newtcolorbox{sectionsummarybox}{enhanced,breakable,colback=CueGray,colframe=RuleGray,boxrule=.55pt,arc=1.5mm,title={Topic Summary},fonttitle=\bfseries\color{Navy},left=2mm,right=2mm,top=1mm,bottom=1mm}
\newtcolorbox{practicebox}{enhanced,breakable,colback=white,colframe=Green,boxrule=.75pt,arc=2mm,title={Practitioner Lens},fonttitle=\bfseries}
\newtcolorbox{scopenotebox}{enhanced,breakable,colback=Warm,colframe=Risk,boxrule=.7pt,arc=2mm,title={Terminology and Scope Note},fonttitle=\bfseries}
\newtcolorbox{synthesisbox}{enhanced,breakable,colback=Navy!5,colframe=Navy,boxrule=.85pt,arc=2mm,title={Chapter Synthesis},fonttitle=\bfseries}
\newtcolorbox{takeawaybox}{enhanced,breakable,colback=Teal!8,colframe=Teal,boxrule=.9pt,arc=2mm,title={One-Sentence Takeaway},fonttitle=\bfseries}
\newcommand{\CornellHeader}{\rowcolor{Navy}\color{white}\bfseries Cue / Active Recall & \color{white}\bfseries Notes / Reasoning \\\endhead}
\newcommand{\CornellRow}[2]{\cellcolor{CueGray}\textbf{#1} & #2 \\\addlinespace[3pt]}

\newcommand{\ChapterMark}{Chapter 4}
\title{Chapter 4: Fundamental Concepts and Models\\[2mm]\large Cornell Notes}
\author{}
\date{}
\hypersetup{pdftitle={Chapter 4: Fundamental Concepts and Models - Cornell Notes},pdfauthor={},pdfsubject={Cloud computing and architecture study notes},pdfkeywords={cloud, architecture, security, operations, Cornell notes}}
\begin{document}
\maketitle
\thispagestyle{fancy}
\vspace{-1.5em}
\begin{orientationbox}
\textbf{Central problem.} Define cloud roles and boundaries, identify characteristic behaviors, and distinguish delivery models from deployment models so responsibility and architecture choices remain precise.\par\smallskip
\textbf{Prerequisites.} Basic cloud terminology, organizational structure, infrastructure layers, application platforms, and service consumption.\par\smallskip
\textbf{Connection.} The models in this chapter determine which mechanisms, controls, operational duties, and cost structures later architectures require.
\end{orientationbox}
\section{Learning Objectives}
\begin{enumerate}
\item Define the principal concepts and decision boundaries in Fundamental Concepts and Models.
\item Explain how 4.1 Roles and boundaries contributes to the chapter's central problem.
\item Distinguish Cloud provider from Cloud consumer using responsibility, behavior, and evidence.
\item Trace components, data, control, identity, capacity, and failure across 4.2 Cloud characteristics.
\item Analyze benefits together with assumptions, limitations, security risks, and operational tradeoffs.
\item Evaluate provider and consumer responsibilities at every relevant trust or administrative boundary.
\item Select verification evidence for architecture, configuration, monitoring, cost, and service-quality claims.
\item Apply the chapter's reasoning to a realistic design, review, migration, or incident scenario.
\end{enumerate}
\section{Cornell Cue-and-Notes}
\Needspace{9\baselineskip}
\subsection{4.1 Roles and boundaries}
\begin{longtable}{@{}C N@{}}
\CornellHeader
\CornellRow{What should be remembered about Cloud provider?}{The provider supplies cloud-based resources or services and operates the supporting layers defined by the service arrangement. \textbf{Review lens.} Verify service scope, administrative powers, subcontractors, control evidence, incident duties, data handling, and termination obligations.}
\CornellRow{What should be remembered about Cloud consumer?}{The consumer contracts for, configures, integrates, or uses cloud services and remains accountable for responsibilities not transferred to the provider. \textbf{Review lens.} Identify consumer-owned identity, data, configuration, workload, monitoring, backup, and governance controls.}
\CornellRow{What should be remembered about Cloud broker?}{A broker mediates between consumers and providers and may aggregate, integrate, negotiate, transform, or manage services. \textbf{Review lens.} Review added trust, credentials, data access, failure dependencies, visibility, contract alignment, and exit paths.}
\CornellRow{What should be remembered about Cloud service owner?}{The service owner is accountable for the cloud service's business purpose, lifecycle, requirements, funding, risk, and performance even when operation is delegated. \textbf{Review lens.} Name one accountable owner and define approval, change, risk acceptance, quality, and retirement decisions.}
\CornellRow{What should be remembered about Cloud resource administrator?}{The resource administrator configures and manages cloud-based IT resources through consoles, APIs, automation, or delegated tools. \textbf{Review lens.} Enforce least privilege, strong authentication, separation of duties, just-in-time access, logging, and break-glass controls.}
\CornellRow{What should be remembered about Organizational boundary?}{An organizational boundary separates resources and responsibilities owned or governed by different organizations or business entities. \textbf{Review lens.} Map contracts, data transfer, support access, legal responsibility, and escalation wherever the boundary is crossed.}
\CornellRow{What should be remembered about Trust boundary?}{A trust boundary exists where data, control, identity, or execution crosses between different assumptions or assurance levels; it may occur inside one organization. \textbf{Review lens.} Place authentication, authorization, validation, encryption, monitoring, and fail-safe handling at each crossing.}
\end{longtable}
\begin{sectionsummarybox}
Roles describe responsibilities within a cloud relationship, while boundaries reveal changes in ownership, administration, policy, and trust.
\end{sectionsummarybox}
\Needspace{9\baselineskip}
\subsection{4.2 Cloud characteristics}
\begin{longtable}{@{}C N@{}}
\CornellHeader
\CornellRow{What should be remembered about On-demand usage?}{Consumers can provision or use resources when needed with limited provider-side manual interaction. \textbf{Review lens.} Review authorization, quotas, policy-as-code, approval exceptions, abuse prevention, and deprovisioning.}
\CornellRow{What should be remembered about Ubiquitous access?}{Services are broadly accessible through standard network mechanisms and multiple consumer device types. \textbf{Review lens.} Account for public exposure, device trust, network variability, authentication, encryption, and location-independent attack paths.}
\CornellRow{What should be remembered about Multitenancy and resource pooling?}{Provider resources serve multiple consumers through logical allocation and isolation, increasing utilization while sharing underlying infrastructure. \textbf{Review lens.} Validate tenant identity, data isolation, noisy-neighbor controls, scheduling, cleanup, encryption, and management-plane separation.}
\CornellRow{What should be remembered about Elasticity?}{Capacity can expand or contract in response to demand, often through automated provisioning and release. \textbf{Review lens.} Test trigger quality, delay, maximums, oscillation, scale-in safety, quota exhaustion, and cost effects.}
\CornellRow{What should be remembered about Measured usage?}{Resource consumption is monitored so it can support reporting, capacity decisions, billing, quotas, or optimization. \textbf{Review lens.} Verify meter accuracy, identity binding, units, aggregation windows, tamper resistance, reconciliation, and dispute handling.}
\CornellRow{What should be remembered about Resiliency?}{The environment can withstand, recover from, or adapt to failures while preserving defined service objectives. \textbf{Review lens.} Confirm independence of replicas, data recovery, switchover, degraded-mode behavior, and exercised restoration.}
\end{longtable}
\begin{sectionsummarybox}
Cloud characteristics describe service behavior expected from the environment; each characteristic creates both value and control requirements.
\end{sectionsummarybox}
\Needspace{9\baselineskip}
\subsection{4.3 Cloud delivery models}
\begin{longtable}{@{}C N@{}}
\CornellHeader
\CornellRow{What should be remembered about Infrastructure as a Service?}{IaaS exposes fundamental compute, storage, and networking resources. The provider manages physical and virtualization layers while the consumer commonly controls operating systems, middleware, applications, identities, and data. \textbf{Review lens.} Document the actual responsibility split because managed add-ons can shift individual duties.}
\CornellRow{What should be remembered about Platform as a Service?}{PaaS provides a managed application platform, runtime, development, deployment, and supporting services. Consumers focus more on application logic, configuration, identities, and data. \textbf{Review lens.} Review platform constraints, supported runtimes, patch responsibility, tenancy, deployment pipeline, observability, portability, and service dependencies.}
\CornellRow{What should be remembered about Software as a Service?}{SaaS delivers a provider-operated application through a service interface. Consumers primarily govern users, configuration, data, integration, and business use. \textbf{Review lens.} Assess identity federation, authorization, tenant configuration, data lifecycle, exports, audit evidence, provider changes, and continuity.}
\CornellRow{What should be remembered about Combined and subordinate models?}{Real solutions layer or combine delivery models, so one service can depend on another provider's platform or infrastructure and responsibility can cross several contracts. \textbf{Review lens.} Trace each dependency to the party that controls it and reconcile downstream guarantees with upstream limitations.}
\end{longtable}
\begin{sectionsummarybox}
Delivery models distinguish the service layer supplied by the provider and therefore redistribute consumer control and responsibility.
\end{sectionsummarybox}
\Needspace{9\baselineskip}
\subsection{4.4 Cloud deployment models}
\begin{longtable}{@{}C N@{}}
\CornellHeader
\CornellRow{What should be remembered about Public cloud?}{A public cloud is owned and operated for use by multiple external consumers, typically through standardized service offerings and shared infrastructure. \textbf{Review lens.} Review tenant isolation, provider governance, internet exposure, region choice, service limits, and contract evidence.}
\CornellRow{What should be remembered about Private cloud?}{A private cloud is dedicated to one organization while retaining cloud characteristics such as self-service, pooling, automation, and measured usage. It may be on or off premises. \textbf{Review lens.} Do not label conventional virtualization a private cloud unless service behavior and operating model support the claim.}
\CornellRow{What should be remembered about Multicloud?}{A multicloud uses services from more than one cloud, often to obtain specialized capabilities, negotiate risk, meet geographic needs, or reduce dependency. \textbf{Review lens.} Evaluate cross-cloud identity, data movement, observability, networking, skills, cost, policy consistency, and operating complexity.}
\CornellRow{What should be remembered about Hybrid cloud?}{A hybrid cloud deliberately integrates cloud and non-cloud or distinct cloud environments so workloads, data, or services cooperate across them. \textbf{Review lens.} Trace cross-environment latency, identity, security, data consistency, failure propagation, monitoring, and SLA alignment.}
\end{longtable}
\begin{sectionsummarybox}
Deployment models describe how cloud environments are owned, accessed, and combined; they do not by themselves determine the delivery model.
\end{sectionsummarybox}
\begin{scopenotebox}
Terminology varies across industry frameworks; these notes retain the chapter's model while requiring explicit definitions in real architecture work.
\end{scopenotebox}
\begin{practicebox}
For each model, mechanism, architecture, metric, or case decision, record the objective, owner, trust boundary, inputs, dependencies, expected behavior, failure consequence, measurement point, response, and retained evidence. A diagram or product label is a hypothesis until runtime and recovery tests support it.
\end{practicebox}
\section{Chapter Recap Matrix}
\begingroup\small\setlength{\tabcolsep}{4pt}
\begin{longtable}{@{}>{\RaggedRight\arraybackslash}p{0.235\textwidth}>{\RaggedRight\arraybackslash}p{0.37\textwidth}>{\RaggedRight\arraybackslash}p{0.325\textwidth}@{}}
\toprule\rowcolor{Navy}\color{white}\textbf{Element} & \color{white}\textbf{Purpose or Behavior} & \color{white}\textbf{Decision and Evidence}\\\midrule\endfirsthead
\toprule\rowcolor{Navy}\color{white}\textbf{Element} & \color{white}\textbf{Purpose or Behavior} & \color{white}\textbf{Decision and Evidence}\\\midrule\endhead
Cloud provider & The provider supplies cloud-based resources or services and operates the supporting layers defined by the service arrangement. & Verify service scope, administrative powers, subcontractors, control evidence, incident duties, data handling, and termination obligations. \\\midrule
Cloud consumer & The consumer contracts for, configures, integrates, or uses cloud services and remains accountable for responsibilities not transferred to the provider. & Identify consumer-owned identity, data, configuration, workload, monitoring, backup, and governance controls. \\\midrule
Cloud broker & A broker mediates between consumers and providers and may aggregate, integrate, negotiate, transform, or manage services. & Review added trust, credentials, data access, failure dependencies, visibility, contract alignment, and exit paths. \\\midrule
Cloud service owner & The service owner is accountable for the cloud service's business purpose, lifecycle, requirements, funding, risk, and performance even when operation is delegated. & Name one accountable owner and define approval, change, risk acceptance, quality, and retirement decisions. \\\midrule
Cloud resource administrator & The resource administrator configures and manages cloud-based IT resources through consoles, APIs, automation, or delegated tools. & Enforce least privilege, strong authentication, separation of duties, just-in-time access, logging, and break-glass controls. \\\midrule
Organizational boundary & An organizational boundary separates resources and responsibilities owned or governed by different organizations or business entities. & Map contracts, data transfer, support access, legal responsibility, and escalation wherever the boundary is crossed. \\\midrule
Trust boundary & A trust boundary exists where data, control, identity, or execution crosses between different assumptions or assurance levels; it may occur inside one organization. & Place authentication, authorization, validation, encryption, monitoring, and fail-safe handling at each crossing. \\\midrule
On-demand usage & Consumers can provision or use resources when needed with limited provider-side manual interaction. & Review authorization, quotas, policy-as-code, approval exceptions, abuse prevention, and deprovisioning. \\\midrule
Ubiquitous access & Services are broadly accessible through standard network mechanisms and multiple consumer device types. & Account for public exposure, device trust, network variability, authentication, encryption, and location-independent attack paths. \\\midrule
Multitenancy and resource pooling & Provider resources serve multiple consumers through logical allocation and isolation, increasing utilization while sharing underlying infrastructure. & Validate tenant identity, data isolation, noisy-neighbor controls, scheduling, cleanup, encryption, and management-plane separation. \\\midrule
Elasticity & Capacity can expand or contract in response to demand, often through automated provisioning and release. & Test trigger quality, delay, maximums, oscillation, scale-in safety, quota exhaustion, and cost effects. \\\midrule
Measured usage & Resource consumption is monitored so it can support reporting, capacity decisions, billing, quotas, or optimization. & Verify meter accuracy, identity binding, units, aggregation windows, tamper resistance, reconciliation, and dispute handling. \\\midrule
Resiliency & The environment can withstand, recover from, or adapt to failures while preserving defined service objectives. & Confirm independence of replicas, data recovery, switchover, degraded-mode behavior, and exercised restoration. \\\midrule
Infrastructure as a Service & IaaS exposes fundamental compute, storage, and networking resources. & Document the actual responsibility split because managed add-ons can shift individual duties. \\\midrule
Platform as a Service & PaaS provides a managed application platform, runtime, development, deployment, and supporting services. & Review platform constraints, supported runtimes, patch responsibility, tenancy, deployment pipeline, observability, portability, and service dependencies. \\\midrule
Software as a Service & SaaS delivers a provider-operated application through a service interface. & Assess identity federation, authorization, tenant configuration, data lifecycle, exports, audit evidence, provider changes, and continuity. \\\midrule
Combined and subordinate models & Real solutions layer or combine delivery models, so one service can depend on another provider's platform or infrastructure and responsibility can cross several contracts. & Trace each dependency to the party that controls it and reconcile downstream guarantees with upstream limitations. \\\midrule
Public cloud & A public cloud is owned and operated for use by multiple external consumers, typically through standardized service offerings and shared infrastructure. & Review tenant isolation, provider governance, internet exposure, region choice, service limits, and contract evidence. \\\midrule
Private cloud & A private cloud is dedicated to one organization while retaining cloud characteristics such as self-service, pooling, automation, and measured usage. & Do not label conventional virtualization a private cloud unless service behavior and operating model support the claim. \\\midrule
Multicloud & A multicloud uses services from more than one cloud, often to obtain specialized capabilities, negotiate risk, meet geographic needs, or reduce dependency. & Evaluate cross-cloud identity, data movement, observability, networking, skills, cost, policy consistency, and operating complexity. \\\midrule
Hybrid cloud & A hybrid cloud deliberately integrates cloud and non-cloud or distinct cloud environments so workloads, data, or services cooperate across them. & Trace cross-environment latency, identity, security, data consistency, failure propagation, monitoring, and SLA alignment. \\\midrule
\bottomrule
\end{longtable}
\endgroup
\section{Practical Review Checklist}
\begin{longtable}{@{}>{\bfseries\color{Teal}}p{0.16\textwidth}>{\RaggedRight\arraybackslash}p{0.77\textwidth}@{}}
\toprule\textbf{Area} & \textbf{Verification action}\\\midrule\endhead
Scope & $\square$ Define the in-scope services, workloads, users, data, environments, regions, and dependencies for Fundamental Concepts and Models. \\\addlinespace[2pt]
Roles & $\square$ Name the provider, consumer, owner, administrator, broker, support, security, finance, and risk responsibilities that apply. \\\addlinespace[2pt]
Boundaries & $\square$ Map organizational, trust, network, tenancy, management, data-location, and failure-domain boundaries. \\\addlinespace[2pt]
Architecture & $\square$ Document how Cloud provider, Cloud consumer, and supporting mechanisms interact through normal, peak, maintenance, and failed states. \\\addlinespace[2pt]
Identity & $\square$ Verify human and workload identities, privileges, credentials, delegation, administrative access, revocation, and audit evidence. \\\addlinespace[2pt]
Data & $\square$ Trace data collection, transfer, processing, storage, replication, backup, retention, export, and deletion. \\\addlinespace[2pt]
Quality & $\square$ Set measurable capacity, latency, availability, reliability, scalability, and recovery criteria relevant to Cloud broker. \\\addlinespace[2pt]
Security & $\square$ Challenge the trust and failure assumptions surrounding Cloud service owner with abuse, bypass, isolation, and compromise scenarios. \\\addlinespace[2pt]
Operations & $\square$ Confirm monitoring, alert ownership, change control, patching, drift management, incident response, and tested recovery. \\\addlinespace[2pt]
Economics & $\square$ Tie resource behavior to meters, units, rates, budgets, anomalies, full-lifecycle cost, and accountable owners. \\\addlinespace[2pt]
Contracts & $\square$ Reconcile architecture behavior with service scope, formulas, exclusions, evidence, notification, remedies, audit rights, and exit terms. \\\addlinespace[2pt]
Verification & $\square$ Preserve reproducible evidence that Cloud resource administrator and the overall design satisfy their intended outcomes under realistic load and failure. \\\addlinespace[2pt]
\bottomrule
\end{longtable}
\begin{synthesisbox}
Cloud models answer different questions. Roles identify who acts; boundaries identify where authority or trust changes; characteristics describe how resources behave; delivery models divide responsibility by service layer; deployment models describe ownership, access, and composition of cloud environments. Mixing these dimensions produces weak security and operating assumptions.
\par\smallskip
\textbf{Cross-chapter application.} The models in this chapter determine which mechanisms, controls, operational duties, and cost structures later architectures require.
\end{synthesisbox}
\section{Self-Test Questions}
\begin{enumerate}
\item What central problem does Fundamental Concepts and Models address?
\item Define Cloud provider and explain why it matters.
\item Compare Cloud consumer with Cloud broker. Which assumptions and responsibilities differ?
\item Trace the normal behavior or decision path associated with Cloud service owner.
\item Which evidence would you request when reviewing Cloud resource administrator?
\item How does this chapter connect to the broader cloud architecture and operating model?
\item Scenario: A team claims that Cloud consumer is implemented but cannot show current configuration, monitoring, or failure-test evidence. What should the reviewer conclude and do next?
\item Scenario: A workload using Cloud broker crosses a provider-consumer boundary and then fails during peak demand. Which responsibilities and failure conditions should be investigated?
\item Why is a named mechanism, model, or architecture insufficient evidence that its intended quality or security outcome is achieved?
\item How would you verify that a change involving Organizational boundary improved the intended outcome without shifting risk or cost elsewhere?
\end{enumerate}
\clearpage
\section{Answer Key}
\begin{enumerate}
\item Define cloud roles and boundaries, identify characteristic behaviors, and distinguish delivery models from deployment models so responsibility and architecture choices remain precise.
\item The provider supplies cloud-based resources or services and operates the supporting layers defined by the service arrangement. The review implication is: Verify service scope, administrative powers, subcontractors, control evidence, incident duties, data handling, and termination obligations.
\item The consumer contracts for, configures, integrates, or uses cloud services and remains accountable for responsibilities not transferred to the provider. In contrast, A broker mediates between consumers and providers and may aggregate, integrate, negotiate, transform, or manage services. Compare control, ownership, state, dependencies, operating evidence, and failure impact rather than names alone.
\item The service owner is accountable for the cloud service's business purpose, lifecycle, requirements, funding, risk, and performance even when operation is delegated. A trace should name initiating actor, input or signal, decision point, participating components, state change, output, monitoring, and failure handling.
\item Useful evidence should demonstrate the review lens: Enforce least privilege, strong authentication, separation of duties, just-in-time access, logging, and break-glass controls. Include architecture, configuration, event, test, and ownership evidence where applicable.
\item The models in this chapter determine which mechanisms, controls, operational duties, and cost structures later architectures require.
\item The implementation claim is unsubstantiated. Request the evidence needed to establish current state, then test positive, negative, overload, and failure paths. For this mechanism specifically: Identify consumer-owned identity, data, configuration, workload, monitoring, backup, and governance controls.
\item Investigate the agreed responsibility split, boundary controls, capacity and quota, workload state, dependent services, provider telemetry, consumer telemetry, escalation, recovery, and whether the contract measures the consumer-visible failure.
\item A label expresses design intent, not runtime proof. Effectiveness depends on placement, configuration, inputs, state, control independence, capacity, monitoring, response, maintenance, and tested behavior during failure.
\item Establish a before-and-after measure, preserve comparable test conditions, verify the expected behavior, inspect adjacent dependencies and costs, and exercise rollback or recovery. Apply this specific lens: Map contracts, data transfer, support access, legal responsibility, and escalation wherever the boundary is crossed.
\end{enumerate}
\section{Key-Term Recap}
\begin{longtable}{@{}>{\RaggedRight\arraybackslash}p{0.27\textwidth}>{\RaggedRight\arraybackslash}p{0.66\textwidth}@{}}
\toprule\rowcolor{Navy}\color{white}\textbf{Term} & \color{white}\textbf{Working meaning}\\\midrule\endhead
Cloud provider & The provider supplies cloud-based resources or services and operates the supporting layers defined by the service arrangement. \\\addlinespace[2pt]
Cloud consumer & The consumer contracts for, configures, integrates, or uses cloud services and remains accountable for responsibilities not transferred to the provider. \\\addlinespace[2pt]
Cloud broker & A broker mediates between consumers and providers and may aggregate, integrate, negotiate, transform, or manage services. \\\addlinespace[2pt]
Cloud service owner & The service owner is accountable for the cloud service's business purpose, lifecycle, requirements, funding, risk, and performance even when operation is delegated. \\\addlinespace[2pt]
Cloud resource administrator & The resource administrator configures and manages cloud-based IT resources through consoles, APIs, automation, or delegated tools. \\\addlinespace[2pt]
Organizational boundary & An organizational boundary separates resources and responsibilities owned or governed by different organizations or business entities. \\\addlinespace[2pt]
Trust boundary & A trust boundary exists where data, control, identity, or execution crosses between different assumptions or assurance levels; it may occur inside one organization. \\\addlinespace[2pt]
On-demand usage & Consumers can provision or use resources when needed with limited provider-side manual interaction. \\\addlinespace[2pt]
Ubiquitous access & Services are broadly accessible through standard network mechanisms and multiple consumer device types. \\\addlinespace[2pt]
Multitenancy and resource pooling & Provider resources serve multiple consumers through logical allocation and isolation, increasing utilization while sharing underlying infrastructure. \\\addlinespace[2pt]
Elasticity & Capacity can expand or contract in response to demand, often through automated provisioning and release. \\\addlinespace[2pt]
Measured usage & Resource consumption is monitored so it can support reporting, capacity decisions, billing, quotas, or optimization. \\\addlinespace[2pt]
Resiliency & The environment can withstand, recover from, or adapt to failures while preserving defined service objectives. \\\addlinespace[2pt]
Infrastructure as a Service & IaaS exposes fundamental compute, storage, and networking resources. The provider manages physical and virtualization layers while the consumer commonly controls operating systems, middleware, applications, identities, and data. \\\addlinespace[2pt]
Platform as a Service & PaaS provides a managed application platform, runtime, development, deployment, and supporting services. Consumers focus more on application logic, configuration, identities, and data. \\\addlinespace[2pt]
Software as a Service & SaaS delivers a provider-operated application through a service interface. Consumers primarily govern users, configuration, data, integration, and business use. \\\addlinespace[2pt]
\bottomrule
\end{longtable}
\begin{takeawaybox}
A sound cloud design names every role, boundary, characteristic, service layer, and deployment relationship before assigning controls or guarantees.
\end{takeawaybox}
\end{document}
